\subsection{Monomials} 
\begin{definition}[Monomial]
A \term{monomial} is a combination of numbers and variables, put together with \emph{multiplication only}. 

In simplest form, we combine all the numeric factors into a single \term{coefficient}, and write repeated variable factors using a single positive integer exponent.

The \term{degree} of a monomial is the total number of variable factors, that is, the sum of all the exponents.
\end{definition}
\begin{example} For each expression, determine whether or not it is a monomial. If it is, write it in simplest form.
\begin{enumerate}\setabc
\item $3x\cdot 5xy\cdot (-6yy)$\makeblanknewf
\item $\frac{2}{3}x^3yxy^5$\makeblanknewf
\item $\dfrac{5y^3}{x}$\makeblanknewf
\item $3x^{-4}y^3$\makeblanknewf
\item $7$\makeblanknewf
\end{enumerate}
\end{example}
\begin{example} For each monomial below, determine the coefficient and degree.
\begin{enumerate}\setabc
\item $4x^2$\lbr Coefficient: \makeblank[0.75in]\hfill Degree: \makeblank[3in]
\item $-\frac{5}{3}x^3y^5$\lbr Coefficient: \makeblank[0.75in]\hfill Degree: \makeblank[3in]
\item $7.4xyzw$\lbr Coefficient: \makeblank[0.75in]\hfill Degree: \makeblank[3in]
\end{enumerate}
\end{example}